\documentclass{beamer}
\usepackage{amsmath, amssymb}
\title{Stochastic Calculus}
\subtitle{Intuition and Derivation of It\^o's Lemma}
\date{}

\begin{document}
%------------------------------------------------
\begin{frame}
\titlepage
\end{frame}

%------------------------------------------------
\begin{frame}{The Intuition: Continuous Jitter}
At the heart of stochastic calculus is the Wiener process (or Brownian motion). If you picture a continuous version of a random walk from a Monte Carlo simulation, that’s exactly what this is.
\medskip

Imagine a dust particle suspended in water, being constantly bombarded by water molecules from every direction. Its path is completely random, unpredictable, and continuous.
\medskip

It doesn't teleport, but every single fraction of a second, it zigs and zags independently of what it just did.
\end{frame}

%------------------------------------------------
\begin{frame}{Why Ordinary Calculus Fails}
\textbf{Intuitive View:} Stochastic paths are continuous but nowhere differentiable. They are ``fractals'' entirely made of sharp corners. No matter how far you zoom in, you cannot draw a tangent line.
\medskip

\textbf{Technical View:} Ordinary calculus (Riemann-Stieltjes integration) requires the integrating variable to have \textit{bounded total variation}. 
\medskip

A Wiener process $W_t$ has \textbf{infinite total variation} on any finite interval, but it has a non-zero, finite \textbf{quadratic variation}:
\[
\lim_{\|\Pi\| \to 0} \sum (W_{t_i} - W_{t_{i-1}})^2 = t
\]
Because the total variation is infinite, ordinary limits of sums fail to converge predictably, breaking standard calculus.
\end{frame}

%------------------------------------------------
\begin{frame}{Ordinary Calculus: Taylor Expansion}
To see exactly how it breaks, we look at the Taylor series. Consider a smooth function $f(x)$ and a small perturbation $\Delta x$:
\[
\Delta f = f(x + \Delta x) - f(x)
\]
Expanding this yields:
\[
\Delta f = f'(x)\Delta x + \frac{1}{2}f''(x)(\Delta x)^2 + \mathcal{O}((\Delta x)^3)
\]
To find the instantaneous rate of change, we divide by $\Delta x$:
\[
\frac{\Delta f}{\Delta x} = f'(x) + \frac{1}{2}f''(x)\Delta x + \mathcal{O}((\Delta x)^2)
\]
\end{frame}

%------------------------------------------------
\begin{frame}{Ordinary Calculus: The Limit}
Taking the limit as $\Delta x \to 0$:
\[
\lim_{\Delta x \to 0} \frac{\Delta f}{\Delta x} = f'(x) + 0 + 0 = f'(x)
\]
Because $(\Delta x)^2$ shrinks infinitely faster than $\Delta x$, all higher-order terms vanish completely.
\medskip

In differential shorthand, we state that $(dx)^2 = 0$, giving us the fundamental rule of ordinary calculus:
\[
df = f'(x) dx
\]
\end{frame}

%------------------------------------------------
\begin{frame}{Stochastic Calculus: Taylor Expansion}
Now consider a function $f(t, X_t)$ where $X_t$ follows a stochastic process defined by a drift $\mu$ and volatility $\sigma$:
\[ dX_t = \mu dt + \sigma dW_t \]
We expand $f$ using a bivariate Taylor series for small changes $\Delta t$ and $\Delta X_t$:
\[
\Delta f = \frac{\partial f}{\partial t}\Delta t + \frac{\partial f}{\partial x}\Delta X_t + \frac{1}{2}\frac{\partial^2 f}{\partial x^2}(\Delta X_t)^2 + \frac{1}{2}\frac{\partial^2 f}{\partial t^2}(\Delta t)^2 + \dots
\]
In ordinary calculus, that $(\Delta X_t)^2$ term would vanish. In stochastic calculus, we must evaluate it carefully.
\end{frame}

%------------------------------------------------
\begin{frame}{The Non-Vanishing Variance}
Let's compute the squared differential $(dX_t)^2$:
\[
(dX_t)^2 = (\mu dt + \sigma dW_t)^2 = \mu^2(dt)^2 + 2\mu\sigma(dt \cdot dW_t) + \sigma^2(dW_t)^2
\]
The random jitter $dW_t$ scales with the square root of time: $dW_t \sim \sqrt{dt}$.
\medskip

Thus, squaring the noise yields a term that scales \textit{linearly} with time:
\[ (dW_t)^2 = dt \]
This gives us the fundamental heuristic rules of stochastic calculus:
\[ (dt)^2 = 0, \qquad dt \cdot dW_t = 0, \qquad (dW_t)^2 = dt \]
\end{frame}

%------------------------------------------------
\begin{frame}{Deriving It\^o's Lemma}
Applying these rules to our squared differential:
\[
(dX_t)^2 = 0 + 0 + \sigma^2 dt
\]
Because $(dW_t)^2 = dt$, the squared term is exactly as ``large'' as the first-order time step $dt$. It survives the limit.
\medskip

Substituting $(dX_t)^2 = \sigma^2 dt$ back into our Taylor expansion, and dropping higher-order terms like $(dt)^2$:
\[
df = \frac{\partial f}{\partial t}dt + \frac{\partial f}{\partial x}dX_t + \frac{1}{2}\frac{\partial^2 f}{\partial x^2}(\sigma^2 dt)
\]
\end{frame}

%------------------------------------------------
\begin{frame}{Where the ``Jitter Penalty'' Comes From}
Because $(dW_t)^2 = dt$, we cannot ignore the second-order terms when we do a Taylor series expansion of our function $f(t, X_t)$.
If we expand $df$ up to the second derivative, we get:
\[ df = \frac{\partial f}{\partial t}dt + \frac{\partial f}{\partial x}dX_t + \frac{1}{2}\frac{\partial^2 f}{\partial x^2}(dX_t)^2 + \dots \]

We substitute our It\^o process definition $dX_t = \mu_t dt + \sigma_t dW_t$ into that $(dX_t)^2$ term:
\begin{align*}
(dX_t)^2 &= (\mu_t dt + \sigma_t dW_t)^2 \\
         &= \mu_t^2(dt)^2 + 2\mu_t\sigma_t(dt)(dW_t) + \sigma_t^2(dW_t)^2
\end{align*}
\end{frame}

%------------------------------------------------
\begin{frame}
In It\^o calculus, $(dt)^2 = 0$ and $dt \cdot dW_t = 0$ because they shrink too fast. But $(dW_t)^2 = dt$. So the equation collapses to:
\[ (dX_t)^2 = \sigma_t^2 dt \]

Substitute this back into the Taylor expansion to get the final theorem:
\[ df(t, X_t) = \left( \frac{\partial f}{\partial t} + \mu_t \frac{\partial f}{\partial x} + \frac{1}{2}\sigma_t^2 \frac{\partial^2 f}{\partial x^2} \right) dt + \sigma_t \frac{\partial f}{\partial x} dW_t \]
\end{frame}

%------------------------------------------------
\begin{frame}{Formal Statement of It\^o's Lemma}
\textbf{Definition (It\^o Process):}
An It\^o process is a stochastic process $X_t$ that can be represented as the integral:
\[ X_t = X_0 + \int_0^t \mu_s ds + \int_0^t \sigma_s dW_s \]
which is commonly written in differential form as $dX_t = \mu_t dt + \sigma_t dW_t$.
\medskip

\textbf{Theorem (It\^o's Lemma):} 
Let $X_t$ be an It\^o process. Let $f(t, x)$ be a scalar function that is twice continuously differentiable in $x$ and once in $t$ (i.e., $f \in C^{1,2}$). Then $Y_t = f(t, X_t)$ is also an It\^o process, and its differential is given by:
\[
df(t, X_t) = \left( \frac{\partial f}{\partial t} + \mu_t \frac{\partial f}{\partial x} + \frac{1}{2}\sigma_t^2 \frac{\partial^2 f}{\partial x^2} \right) dt + \sigma_t \frac{\partial f}{\partial x} dW_t
\]
\medskip
\end{frame}
%------------------------------------------------
%------------------------------------------------
\begin{frame}{The Model: Geometric Brownian Motion (GBM)}
We want to model a quantity $S_t$ (like a stock price or population) that grows at a constant average rate but is subject to continuous random shocks proportional to its size.

\vspace{0.5cm}
\textbf{The Stochastic Differential Equation (SDE):}
\[ dS_t = \mu S_t dt + \sigma S_t dW_t \]

\textbf{Formal Definitions:}
\begin{itemize}
    \item $dS_t$: The instantaneous change in the variable $S_t$.
    \item $\mu$: The constant percentage drift (expected growth rate).
    \item $\sigma$: The constant percentage volatility (intensity of the noise).
    \item $dW_t$: The increment of a standard Wiener process (the random shock).
\end{itemize}
\end{frame}

%------------------------------------------------
\begin{frame}{The Transformation Strategy}
Because $S_t$ appears on both the left side and in the differential terms on the right, we cannot integrate directly. We must isolate $S_t$.

\vspace{0.3cm}
\textbf{The Strategy:} Apply a transformation using the natural logarithm, $f(S_t) = \ln(S_t)$. 

\vspace{0.3cm}
First, we compute the necessary partial derivatives of $f(S) = \ln(S)$:
\[
\frac{\partial f}{\partial S} = \frac{1}{S}, \qquad \frac{\partial^2 f}{\partial S^2} = -\frac{1}{S^2}
\]

\vspace{0.3cm}
Let's see what happens when we apply the chain rule in both ordinary and stochastic calculus.
\end{frame}

%------------------------------------------------
\begin{frame}{Side-by-Side: The System Setup}
\begin{columns}[T]
\begin{column}{0.45\textwidth}
\textbf{Ordinary Calculus} \\
(Naive Approach)
\vspace{0.3cm}

Apply ordinary calculus rules. The system is governed by the SDE:
\[ dS_t = \mu S_t dt + \sigma S_t dW_t \]
\end{column}
\hfill\vrule\hfill
\begin{column}{0.45\textwidth}
\textbf{Stochastic Calculus} \\
(Stochastic System)
\vspace{0.3cm}

Volatility is non-zero ($\sigma > 0$). The system is governed by an SDE:
\[ dS_t = \mu S_t dt + \sigma S_t dW_t \]
\end{column}
\end{columns}
\end{frame}

%------------------------------------------------
\begin{frame}{Side-by-Side: Applying the Chain Rule}
\begin{columns}[T]
\begin{column}{0.45\textwidth}
\textbf{Ordinary Calculus}
\vspace{0.3cm}

Using the standard first-order chain rule for $f(S_t) = \ln(S_t)$:
\[ d(\ln S_t) = \frac{\partial f}{\partial S} dS_t \]

Substitute the derivative:
\[ d(\ln S_t) = \frac{1}{S_t} dS_t \]
\end{column}
\hfill\vrule\hfill
\begin{column}{0.45\textwidth}
\textbf{Stochastic Calculus}
\vspace{0.3cm}

Using It\^o's Lemma (second-order expansion) for $f(S_t) = \ln(S_t)$:
\[ d(\ln S_t) = \frac{\partial f}{\partial S} dS_t + \frac{1}{2}\frac{\partial^2 f}{\partial S^2}(dS_t)^2 \]

Substitute the derivatives:
\[ d(\ln S_t) = \frac{1}{S_t} dS_t - \frac{1}{2S_t^2}(dS_t)^2 \]
\end{column}
\end{columns}
\end{frame}

%------------------------------------------------
\begin{frame}{Evaluating the Quadratic Term}
In ordinary calculus, $(dS_t)^2 = 0$, so the second term vanishes. In stochastic calculus, we must evaluate $(dS_t)^2$ explicitly.

\vspace{0.3cm}
Square the full SDE:
\[ (dS_t)^2 = (\mu S_t dt + \sigma S_t dW_t)^2 \]
\[ (dS_t)^2 = \mu^2 S_t^2 (dt)^2 + 2\mu\sigma S_t^2 (dt \cdot dW_t) + \sigma^2 S_t^2 (dW_t)^2 \]

\vspace{0.3cm}
Apply the fundamental rules: $(dt)^2 = 0$, $dt \cdot dW_t = 0$, and crucially, $(dW_t)^2 = dt$.
\[ (dS_t)^2 = 0 + 0 + \sigma^2 S_t^2 dt \]
\[ (dS_t)^2 = \sigma^2 S_t^2 dt \]
\end{frame}

%------------------------------------------------
\begin{frame}{Side-by-Side: Substituting $dS_t$ Back In}
\begin{columns}[T]
\begin{column}{0.45\textwidth}
\textbf{Ordinary Calculus}
\vspace{0.3cm}

\[ d(\ln S_t) = \frac{1}{S_t} (\mu S_t dt + \sigma S_t dW_t) \]

The $S_t$ variables cancel out:
\[ d(\ln S_t) = \mu dt + \sigma dW_t \]
\end{column}
\hfill\vrule\hfill
\begin{column}{0.45\textwidth}
\textbf{Stochastic Calculus}
\vspace{0.3cm}

\[ d(\ln S_t) = \frac{1}{S_t} dS_t - \frac{1}{2S_t^2} (\sigma^2 S_t^2 dt) \]

Substitute $dS_t$:
\begin{align*}
d(\ln S_t) &= \frac{1}{S_t}(\mu S_t dt + \sigma S_t dW_t) \\
&\quad - \frac{1}{2}\sigma^2 dt
\end{align*}

The $S_t$ variables cancel out:
\[ d(\ln S_t) = \mu dt + \sigma dW_t - \frac{1}{2}\sigma^2 dt \]
\end{column}
\end{columns}
\end{frame}

%------------------------------------------------
\begin{frame}{Side-by-Side: Integration}
\begin{columns}[T]
\begin{column}{0.45\textwidth}
\textbf{Ordinary Calculus}
\vspace{0.3cm}

The variable $S_t$ is now completely isolated. We integrate both sides from $0$ to $t$:
\[ \int_0^t d(\ln S_u) = \int_0^t (\mu du + \sigma dW_u) \]
\[ \ln(S_t) - \ln(S_0) = \mu t + \sigma W_t \]
\end{column}
\hfill\vrule\hfill
\begin{column}{0.45\textwidth}
\textbf{Stochastic Calculus}
\vspace{0.3cm}

The variable $S_t$ is now completely isolated. We integrate both sides from $0$ to $t$:
\[ \int_0^t d(\ln S_u) = \int_0^t \left(\mu - \frac{1}{2}\sigma^2\right) du 

\hfill + \int_0^t \sigma dW_u \]
\[ \ln(S_t) - \ln(S_0) = \left(\mu - \frac{1}{2}\sigma^2\right)t 

\hfill \hfill + \sigma W_t \]
\end{column}
\end{columns}
\end{frame}

%------------------------------------------------
\begin{frame}{Side-by-Side: The Final Explicit Solutions}
\begin{columns}[T]
\begin{column}{0.45\textwidth}
\textbf{Ordinary Calculus}
\vspace{0.3cm}

Exponentiate both sides to find $S_t$:
\[ S_t = S_0 \exp(\mu t + \sigma W_t) \]

\vspace{1.2cm}
\textit{Result:} Mathematically incorrect naive guess. Ordinary calculus wrongly assumes the continuous variance carries no penalty.
\end{column}
\hfill\vrule\hfill
\begin{column}{0.45\textwidth}
\textbf{Stochastic Calculus}
\vspace{0.3cm}

Exponentiate both sides to find $S_t$:
\[ S_t = S_0 \exp\left( \left(\mu - \frac{1}{2}\sigma^2\right)t 

\hfill + \sigma W_t \right) \]

\vspace{0.3cm}
\textit{Result:} A log-normal process. The continuous noise inherently degrades the growth trajectory.
\end{column}
\end{columns}
\end{frame}

%------------------------------------------------
\begin{frame}{The Conclusion: Volatility Drag}
By solving the models side-by-side, we see exactly what ordinary calculus missed. 

\vspace{0.5cm}
If we simply added random noise to the ordinary calculus solution, we would have guessed $S_t = S_0 \exp(\mu t + \sigma W_t)$. \textbf{This is mathematically incorrect.}

\vspace{0.5cm}
The rigorous application of It\^o's Lemma reveals the $-\frac{1}{2}\sigma^2$ term. This is known as the \textbf{volatility drag}. It proves that the geometric mean growth rate of a stochastic process is strictly less than its arithmetic mean drift $\mu$, entirely due to the accumulated penalty of the path's non-zero quadratic variation.
\end{frame}

\end{document}